\section{Introduction}
\label{sec:introduction}

Let \(V\) be a finite-dimensional vector space over \(\F_2\). We consider
the functional equation
\begin{equation}\label{eq:functional}
 f(f(x))=f(x)+x\qquad(x\in V)
\end{equation}
under the assumption that \(f\) has binary algebraic degree at most two.
Every nonzero orbit is the triple \(\{x,f(x),x+f(x)\}\), the nonzero
points of a binary plane. The equation therefore ties a local cycle
condition to the additive geometry of the whole space. We ask how a
quadratic degree bound organizes these orbits into an algebraic structure
that can be classified and measured.

The linear case identifies the relevant scalar field. A linear solution
has minimal polynomial \(X^2+X+1\) on every nonzero space, so it defines
an intrinsic \(\F_4\)-structure in which \(f\) is multiplication by
\(\omega\), where \(\omega^2+\omega=1\). Up to binary-linear conjugacy,
there is one such solution in each positive even dimension.
Degree two is the first level at which the polar form can be nonzero.
It introduces interactions between the orbit planes, and these
interactions are the subject of the paper.

The first structural object is the polar form
\[
 [x,y]=f(x+y)+f(x)+f(y).
\]
For a general zero-constant quadratic map this form is alternating and
bilinear. Here the functional equation also forces Jacobi, making it a
Lie bracket. The derived space, the center, and the lower central series
therefore become intrinsic invariants of the map. The equation imposes
further rigidity: solvability implies nilpotency, and fourth polarization
gives a degree-four polynomial identity. This polar structure is the
bridge from the permutation equation to both classification and sharp
dimension bounds.

\subsection{The main correspondence}

Write \(W=[V,V]\) for the binary span of all polar values.
The principal classification theorem concerns the metabelian case
\([W,W]=0\). The maps induced by \(f\) on \(W\) and \(V/W\) are linear
operators satisfying \(T^2+T=\id\); they therefore define intrinsic
\(\F_4\)-vector space structures. Fix
\(\omega\in\F_4\) with \(\omega^2+\omega=1\).

In suitable binary-linear coordinates, every metabelian solution has the form
\begin{equation}\label{eq:intro-normal}
 f(x,z)=\bigl(\omega x,\,
       \omega z+P(x)z+Q(x)+P(x)Q(x)\bigr),
\end{equation}
where
\[
 P(x)=\sum_{i=1}^d x_iP_i,
 \qquad Q(x)=\sum_{i<j}x_ix_ja_{ij}.
\]
The \(\F_4\)-linear operators \(P_i\) commute, satisfy \(P_i^2=0\), and obey
\[
 P_i a_{jk}+P_j a_{ik}+P_k a_{ij}=0.
\]
The coefficients \(a_{ij}\) generate \(W\) under the operators \(P_i\).
Conversely, all such data give solutions.

The key necessity statement is that binary degree two forces the
square-zero actions on the intrinsic polar image \(W\). Thus the normal
form describes every metabelian solution as well as constructing them.
Moreover, the correspondence includes morphisms. Binary-linear conjugacy
of these maps is equivalent to \(\F_4\)-Lie isomorphism of the associated
metabelian algebras with \((\ad u)^2|_{W}=0\).
In particular it identifies the linear centralizer of \(f\) with the
automorphism group of the associated \(\F_4\)-Lie algebra
(Theorem~\ref{thm:correspondence}).

\subsection{Sharp dimension costs}

The normal form gives an exact measure of the dimension required for a
given polar Lie nilpotency class. Let \(\mu(c)\) be the minimum binary
dimension among metabelian solutions of class \(c\ge2\). We prove
\begin{equation}\label{eq:intro-minimum}
 \mu(2)=6,\qquad\mu(3)=8,\qquad\mu(4)=12,\qquad
 \mu(c)=2^c+2c-6\quad(c\ge5).
\end{equation}
The proof uses a universal commutator module over a square-zero polynomial
ring. Its multigrading determines the possible terminal commutators, and a
dual contraction argument determines the smallest quotient that retains
one of them. Both the lower bound and its attainment are uniform in \(c\)
(Theorem~\ref{thm:unrestricted-minimum}).

The \(\F_4\)-Lie algebra in this correspondence is auxiliary: the original
binary polar bracket need not itself be \(\F_4\)-bilinear. We give an exact
linear criterion for when it admits such a structure with the prescribed
intrinsic scalar actions on \(W\) and \(V/W\).
This condition is equivalent to the existence of a restricted square on
the auxiliary algebra that vanishes on \(W\) and takes values in \(W\)
(Theorem~\ref{thm:compatibility}).

Let \(\nu(c)\) denote the minimum dimension in this compatible subclass.
We obtain the single formula
\begin{equation}\label{eq:intro-compatible-minimum}
 \nu(c)=2^c+2c-2\qquad(c\ge2).
\end{equation}
One part of its lower bound counts independent preimages of the nonempty
faces of a square-zero operator cube. This gives
\(\dim_{\F_4}W\ge2^{c-1}-1\), independently of the unrestricted extremal
argument. The other part proves \(\dim_{\F_4}(V/W)\ge c\).
The elementary augmentation-module construction attains both bounds
(Theorem~\ref{thm:compatible-minimum}).

\begin{table}[ht]
\centering
\caption{Minimum binary dimensions in the two metabelian subclasses.}
\label{tab:minima}
\begin{tabular}{@{}lrrrrrr@{}}
\toprule
Polar class \(c\)&2&3&4&5&6&7\\
\midrule
\(\mu(c)\)&6&8&12&36&70&136\\
\(\nu(c)\)&6&12&22&40&74&140\\
\bottomrule
\end{tabular}
\end{table}

Compatibility therefore has an exact additional cost of four binary
dimensions in every class \(c\ge5\). Equality still allows growing
structural diversity: for each \(c\ge4\) we construct at least \(c-2\)
pairwise inequivalent compatible extremizers
(Theorem~\ref{thm:extremizers}).

\subsection{Closed classification ranges}

Two applications illustrate the reach and the boundaries of the theory.
There are exactly three binary-linear conjugacy classes in dimension eight
(Theorem~\ref{thm:dimension-eight}). The proof combines the correspondence
with the classical small-dimensional simple Lie algebra lists and short
explicit obstructions to the polar identities.

There is also a closed classification in an all-dimensional stratum:
if the polar algebra is solvable and
\(\dim_{\F_2}(V/W)=4\), exactly one class occurs in each of dimensions
\(4,6,8,10,12,14,16\), and none occurs thereafter
(Theorem~\ref{thm:seven-chain}). Its final two members are nonmetabelian.
Completeness is proved by a lossless central descent, exact extension-space
calculations, and explicit projective orbit actions. The finite rank portion
is supplied with independently replayable certificates.

\subsection{Context and related work}

There is an operator viewpoint on \eqref{eq:functional}. On complete
mappings, \(\mathcal F(g)=(\id+g)^{-1}\) satisfies
\(\mathcal F^3(g)=g\), and its fixed objects are exactly the solutions
of \eqref{eq:functional}. The arithmetic study of norm fibers and affine
companion groups in~\cite{Li2026} also describes these fixed objects by
oriented binary line spreads. The present paper studies their quadratic
structure through the polar Lie algebra.

Equation~\eqref{eq:functional} implies that both \(f\) and \(f+\id\) are
permutations. Triple-cycle and regular complete permutation polynomials
have been studied by Liu--Chen--Xu--Sun~\cite{Liu2019},
Xu--Zeng--Zhang~\cite{Xu2022}, and Lu--Wu--Wang--Cao~\cite{Lu2023}.
Quadratic complete mappings also occur in the constructions of
Samardjiska--Gligoroski~\cite{Samardjiska2017}.
Here the orbit-sum identity, binary degree, and binary-linear conjugacy
are considered together. In particular, simultaneous three-cycle behavior
of \(f\) and \(f+\id\) does not imply \eqref{eq:functional};
Example~\ref{ex:regular-cpp} gives a quadratic instance from the
affine-switching construction of~\cite{Xu2022}.

There is a parallel quasigroup interpretation. The operation
\(x*y=x+f(x+y)\) is a quaternary Mendelsohn quasigroup.
Zhang's thesis~\cite[Chapter~3]{Zhang1998} develops central representations
and recursions in that setting. Its Lemma~10.3(1) already contains the
central six-dimensional example used below. Throughout this paper,
equivalence means binary-linear conjugacy, and nilpotency class means
the class of the polar Lie algebra. These specifications retain the
addition and the degree filtration that enter the necessity proofs.

The commutator-module presentation is the classical Koszul presentation
of the free metabelian derived module; see
Papadima--Suciu~\cite[Theorem~6.2]{PapadimaSuciu2004} and the
characteristic-two module methods of Kov\'acs--St\"ohr~\cite{KovacsStohr2013}.
Restricted squares and quadratic maps have a broader established
connection; see, for example, Pakianathan--Yal\c{c}{\i}n~\cite{Pakianathan2007}.
We use these tools to prove the degree-forced correspondence and the exact
extremal formulas for \eqref{eq:functional}. The four-variable polar identity
is the classical degree-four standard polynomial in Lie form
\cite[Theorem~9]{Drensky2021}; its forced occurrence here provides the
obstructions used in the low-dimensional application.

\subsection{Conventions}

We use \(f^j\) only for composition powers. All dimensions without a
specified field are binary dimensions. A map is called quadratic when its
binary algebraic degree is at most two, so the linear cases are included.
Two solutions are equivalent if \(Sf=f'S\) for a binary-linear isomorphism
\(S\). Lie brackets are alternating; in characteristic two they are
symmetric as bilinear maps. We put
\(\gamma_1(L)=L\), \(\gamma_{i+1}(L)=[L,\gamma_i(L)]\), and
\(L^{(0)}=L\), \(L^{(i+1)}=[L^{(i)},L^{(i)}]\).
The class of a nonzero abelian algebra is one.

\medskip
\noindent\textbf{Notation at a glance.}
The \(\F_4\)-dimensions below are used in the metabelian correspondence.
\begin{center}
\small
\renewcommand{\arraystretch}{1.12}
\begin{tabular}{@{}p{.25\textwidth}p{.71\textwidth}@{}}
\toprule
Symbol & Meaning\\
\midrule
\(L=(V,[\ ,\ ])\) & The binary polar Lie algebra of \(f\).\\
\(W=L',\ U=V/W\) & The intrinsic derived ideal and abelianization.\\
\(d,\ w\) & \(\dim_{\F_4}U\) and \(\dim_{\F_4}W\).\\
\(\omega\) & A root of \(X^2+X+1\) in \(\F_4\).\\
\(\m\) & The augmentation ideal \((t_1,\ldots,t_d)\) of \(R_d\),
or its image under the operator action.\\
\(\gamma_i(L),\ L^{(i)}\) & The lower central and derived series.\\
\(\mu(c),\ \nu(c)\) & Minimum binary dimensions at class \(c\):
metabelian and compatible metabelian, respectively.\\
\(\mathcal A_k\) & Filtered action spaces used to distinguish extremizers
in Theorem~\ref{thm:extremizers}.\\
\bottomrule
\end{tabular}
\end{center}
